\Needspace{12\baselineskip}\SourceLeaf{315}

\begin{center}
[2.]
\end{center}

\begin{center}
\textsc{Gauss} an \textsc{Gerling}. Göttingen, 7. November 1830.
\end{center}

\indent\mbox{. . . . . . . . . .} Über die vier merkwürdigen Punkte im Dreieck
habe ich selbst nie etwas bekannt gemacht. Was ich jedoch 1808 darüber
\textsc{Schumacher} mitgetheilt habe, hat dieser im Anhange zu seiner
Übersetzung von \textsc{Carnot}, Géométrie de position, zweiter
Bd., welcher 1810 erschienen ist, drucken lassen. Ich zweifle fast, dass
der Punkt, welcher sich nach dem von Ihnen angeführten Satze ergibt,
besonders elegante Relationen zu den übrigen darbieten wird, und zwar
deshalb, weil jenes Theorem nur ein ganz specieller Fall eines viel
allgemeinern ist.

Wenn man nemlich an die Seiten eines gegebenen Dreiecks $ABC$
drei andere $AB\gamma$, $\alpha BC$, $AC\beta$ legt (entweder alle drei
an die äussere oder alle drei an die innere Seite des Dreiecks
$ABC$ gelegt), die jedes einem zweiten gegebenen Dreieck $abc$ ähnlich
sind, und zwar dergestalt, dass sie auch ähnlich liegend und die Winkel an
$\alpha$, $\beta$, $\gamma$ denen an $a$, $b$, $c$ gleich werden, so
schneiden sich die Linien $A\alpha$, $B\beta$, $C\gamma$
(nöthigenfalls vorwärts oder rückwärts fortgesetzt) in Einem Punkte
$D$ und die Winkel an $D$ ($ADB$, $BDC$, $CDA$) sind
denen an $\gamma$, $\alpha$, $\beta$ entweder selbst oder ihren
Nebenwinkeln gleich und ähnlich liegend.

Offenbar ist diess die zierlichste Construction der
\textsc{Pothenots}chen Aufgabe; wo man aber in der Praxis sich begnügen
kann, nur Einen Punkt z.\,B. $\alpha$ zu construiren und den Messtisch nach
der Linie $A\alpha$ zu orientiren. Wie ich von meinem Sohne höre, ist
letztere Manier unter dem Namen der \textsc{Schulz-Montanus}'schen
bekannt. \mbox{. . . . . . .}
